\documentclass[12pt,a4paper]{article}
\usepackage[utf8]{inputenc}
\usepackage[T1]{fontenc}
\usepackage{lmodern}
\usepackage[french]{babel}
\usepackage[margin=2.2cm]{geometry}
\usepackage{amsmath,amssymb}
\usepackage{enumitem}
\usepackage[dvipsnames]{xcolor}
\usepackage[most]{tcolorbox}
\setlength{\parskip}{0.5em}\setlength{\parindent}{0pt}
\setlist{nosep}
\pagestyle{empty}

\newcommand{\Var}{\operatorname{Var}}\newcommand{\Cov}{\operatorname{Cov}}
\newtcolorbox{outil}[1][]{colback=ForestGreen!6,colframe=ForestGreen!70!black,fonttitle=\bfseries,#1}

\begin{document}

\begin{center}
{\LARGE\bfseries Chapitre 3 -- Tests d'hypothèse}\\[4pt]
{\large La fiche simple}
\end{center}

\section*{1. L'idée}

On a estimé un coefficient, par exemple $\hat b = 0{,}2$. On veut savoir si le vrai coefficient pourrait être 0.

\textbf{On se pose une seule question :} si le vrai coefficient était 0, est-ce qu'obtenir 0,2 serait \emph{rare} ?
\begin{itemize}
\item Si c'est rare (moins de 5\,\% de chances) : on \textbf{rejette} $H_0$.
\item Sinon : on \textbf{ne rejette pas} $H_0$.
\end{itemize}

\section*{2. Les outils}

\begin{outil}[title={Test $t$ : tester UN coefficient}]
\textbf{Formule :}\quad $t=\dfrac{\text{estimation}-\text{valeur testée}}{\text{écart-type}}=\dfrac{b_k-\beta_{k,0}}{\text{e.t.}(b_k)}$

\textbf{Règle :}\quad on rejette $H_0$ si $|t|>$ valeur critique \quad (environ \textbf{1,96} à 5\,\%).
\end{outil}

\begin{outil}[title={$p$-value : la même décision, autrement}]
\textbf{En mots :} la probabilité d'obtenir un résultat aussi extrême si $H_0$ est vraie.

\textbf{Règle :}\quad on rejette $H_0$ si $p<0{,}05$.
\end{outil}

\begin{outil}[title={Intervalle de confiance à 95\,\%}]
\textbf{Formule :}\quad $b_k\ \pm\ 1{,}96\times\text{e.t.}(b_k)$

\textbf{Règle :}\quad si 0 n'est \textbf{pas} dans l'intervalle, on rejette $\beta_k=0$.
\end{outil}

\begin{outil}[title={Tester une SOMME de coefficients : $\beta_1+\beta_2=1$}]
\textbf{Formule :}\quad $t=\dfrac{b_1+b_2-1}{\text{e.t.}(b_1+b_2)}$, \quad avec \quad $\Var(b_1+b_2)=\Var(b_1)+\Var(b_2)+\mathbf{2\,Cov}(b_1,b_2)$

\textbf{À ne pas oublier :} la covariance.
\end{outil}

\begin{outil}[title={Test $F$ : tester PLUSIEURS restrictions à la fois}]
\textbf{Formule :}\quad $F=\dfrac{\text{SCE}_r-\text{SCE}_u}{\text{SCE}_u}\times\dfrac{n-K}{J}$

$\text{SCE}_r$ : modèle avec les restrictions ; $\text{SCE}_u$ : modèle sans restriction ; $J$ : nombre de restrictions.

\textbf{Règle :}\quad on rejette $H_0$ si $F>$ valeur critique.
\end{outil}

\newpage
\section*{3. Un exemple (celui du cours)}

$H_0:\beta_3=0{,}5$. On a trouvé $b_3=0{,}77$ et $\text{e.t.}(b_3)=0{,}12$, avec $n-K=11$.

\begin{enumerate}
\item $t=\dfrac{0{,}77-0{,}5}{0{,}12}=2{,}25$
\item Valeur critique à 5\,\% avec 11 degrés de liberté : 2,201.
\item $2{,}25>2{,}201$ : on \textbf{rejette} $H_0$. (La $p$-value vaut 0,046, qui est bien inférieure à 0,05.)
\end{enumerate}

\section*{4. Les 5 choses à retenir}

\begin{enumerate}[itemsep=4pt]
\item \textbf{Erreur de type I} = rejeter $H_0$ alors qu'elle est vraie. Elle arrive avec une probabilité de 5\,\% si $\alpha=5\,\%$.
\item \textbf{Degrés de liberté} = $n-K$ (nombre d'observations moins nombre de coefficients).
\item \textbf{Grand échantillon} : on utilise 1,96 à 5\,\%. \textbf{Petit échantillon} : on lit la table de Student avec $n-K$.
\item \textbf{Intervalle de confiance} : on ne dit pas « 95\,\% de chances que $\beta$ soit dedans », mais « 95\,\% des intervalles construits ainsi contiennent le vrai $\beta$ ».
\item \textbf{Variables très corrélées} (multicolinéarité) : les tests $t$ peuvent être non significatifs alors que le test $F$ l'est. Cela veut dire qu'au moins une variable compte, sans savoir laquelle.
\end{enumerate}

\section*{5. Les compléments utiles à l'examen}

\begin{outil}[title={Reparamétrisation : faire le test d'une somme avec un simple $t$}]
\textbf{À quoi ça sert :} tester $\beta_1+\beta_2=1$ sans calculer la covariance.

\textbf{Méthode :} poser $\theta=\beta_1+\beta_2-1$, donc $\beta_1=\theta-\beta_2+1$, et remplacer dans le modèle :
\[
y_i-x_{1i}=\beta_0+\theta\,x_{1i}+\beta_2\,(x_{2i}-x_{1i})+\epsilon_i
\]
\textbf{Règle :} le $t$ du coefficient de $x_{1i}$ teste directement $\theta=0$, c'est-à-dire $\beta_1+\beta_2=1$.
\end{outil}

\begin{outil}[title={Test unilatéral : quand l'alternative a un seul sens}]
\textbf{Exemple :} $H_0:\beta_k\le0$ contre $H_1:\beta_k>0$.

\textbf{Règle :} on rejette si $t>1{,}65$ (à 5\,\%, un seul côté). La $p$-value est la moitié de la $p$-value bilatérale.
\end{outil}

\begin{outil}[title={Test $F$ : toutes les pentes sont-elles nulles ?}]
\textbf{Hypothèse :} $H_0$ : tous les coefficients sauf la constante valent 0.

\textbf{Formule :} même $F$ que page 1, avec $J=K-1$ ; le modèle contraint est $y$ sur une constante seulement.

\textbf{Règle :} rejet si $F>F^*(K-1,\,n-K)$. Les logiciels l'affichent en haut du tableau (« F » et « Prob $>$ F »).
\end{outil}

\begin{outil}[title={Écrire plusieurs restrictions sous la forme $R\beta=q$}]
\textbf{Exemple :} $H_0:\beta_2=1$ et $\beta_3=1$ dans $y=\beta_0+\beta_1x_1+\beta_2x_2+\beta_3x_3+\epsilon$ :
\[
\underbrace{\begin{pmatrix}0&0&1&0\\0&0&0&1\end{pmatrix}}_{R}\begin{pmatrix}\beta_0\\\beta_1\\\beta_2\\\beta_3\end{pmatrix}=\underbrace{\begin{pmatrix}1\\1\end{pmatrix}}_{q}
\]
\textbf{Règle :} $R$ a une ligne par restriction, donc $J$ = nombre de lignes. On transforme chaque restriction en « un coefficient vaut 0 » (comme la reparamétrisation), puis on fait le test $F$.
\end{outil}

\begin{outil}[title={Multicolinéarité : pourquoi la variance augmente}]
\textbf{Formule :}\quad $\Var(b_k)=\dfrac{\sigma^2}{x_k'x_k\,(1-R_k^2)}$

$R_k^2$ : $R^2$ (non centré) de la régression de $x_k$ sur les autres variables.

\textbf{En mots :} plus $x_k$ ressemble aux autres variables ($R_k^2$ proche de 1), plus sa variance est grande et plus son $t$ est petit.
\end{outil}

\begin{outil}[title={Après un Monte Carlo : l'estimateur est-il sans biais ?}]
\textbf{Formule :}\quad $\tau=\dfrac{\bar b^*-\beta_k}{\sqrt{\widehat{\Var}(b^*)/R}}$ \quad ($\bar b^*$ : moyenne des $R$ estimations simulées)

\textbf{Règle :} si $|\tau|<1{,}96$, on ne rejette pas l'absence de biais.
\end{outil}

\end{document}
